\section{Анализ предметной области и технологического стека}

\subsection{Исходные данные}

Работа выполняется на открытом датасете взаимодействий пользователей с короткими видео.
Полный объем исходных данных составляет около 1.5~ТБ (десятки миллиардов событий, миллионы пользователей и видео) --- обработка такого объема на одной машине нецелесообразна, поэтому в соответствии с заданием практики использована репрезентативная часть датасета ($\approx$2.6~ГБ).
Прореживание выполнено по пользователям и видео с сохранением всех событий отобранных сущностей, поэтому распределения активности и плотность взаимодействий в выборке соответствуют полным данным, а пайплайн спроектирован потоковым и масштабируется на бóльшую долю датасета изменением одного параметра конфигурации.

Используемая выборка охватывает 27 последовательных недель и содержит:

\begin{itemize}
    \item \textbf{47\,976\,280 событий} взаимодействия (10\,000 пользователей, 19\,628 видео);
    \item по каждому событию --- время просмотра \texttt{timespent}, явный фидбек (\texttt{like}, \texttt{dislike}, \texttt{share}, \texttt{bookmark}, \texttt{click\_on\_author}, \texttt{open\_comments}) и контекст показа (\texttt{place}, \texttt{platform}, \texttt{agent});
    \item метаданные пользователей (возраст, пол, регион) и видео (длительность, идентификатор автора);
    \item контентные эмбеддинги видео размерности 64 в формате \texttt{npz}; в работе используются первые 32 компоненты (Matryoshka-представление).
\end{itemize}

Ключевое свойство датасета --- глобальный временнóй порядок: события отсортированы по времени как внутри недельных шардов, так и между ними.
Это позволяет строить темпоральные разбиения train/validation/test и признаки без утечки будущего.

Структура исходных файлов приведена в таблице~\ref{tab:raw-files}.

\begin{table}[H]
\centering
\caption{Структура исходных данных}
\label{tab:raw-files}
\small
\begin{tabularx}{\textwidth}{@{}l X@{}}
\toprule
\textbf{Файлы} & \textbf{Содержимое} \\
\midrule
\texttt{**/week\_NN.parquet} & Недельные шарды взаимодействий: \texttt{user\_id}, \texttt{item\_id}, \texttt{timespent}, фидбек, контекст \\
\texttt{metadata/users\_metadata.parquet} & Пользователи: \texttt{age}, \texttt{gender}, \texttt{geo} \\
\texttt{metadata/items\_metadata.parquet} & Видео: \texttt{duration}, \texttt{author\_id} \\
\texttt{metadata/item\_embeddings.npz} & Массивы \texttt{item\_id} и \texttt{embedding} (контентные эмбеддинги) \\
\bottomrule
\end{tabularx}
\end{table}

Особенность хранения --- компактные типы: идентификаторы \texttt{UInt32}, время просмотра и длительность \texttt{UInt8}, фидбек \texttt{Boolean}, категориальные атрибуты --- целочисленные коды \texttt{UInt8}.
Это в разы сокращает объем на диске и в памяти по сравнению со строковыми представлениями.

\subsection{Задача ранжирования}

Задача формулируется как learning-to-rank: для каждого пользователя (группы) модель должна упорядочить показанные ему видео так, чтобы более релевантные оказались выше.
В отличие от бинарной классификации, оптимизируется именно порядок внутри группы, а основной метрикой качества служит NDCG@k.

Целевая переменная (релевантность) строится из неявного фидбека: досмотр видео и позитивные действия увеличивают релевантность с настраиваемыми весами, дизлайк обнуляет ее.
Такой градуированный таргет информативнее бинарного <<клик/не клик>> и естественно сочетается с listwise-лоссами.

\subsection{Технологический стек}

Выбранный стек приведен в таблице~\ref{tab:stack}.

\begin{table}[H]
\centering
\caption{Технологический стек решения}
\label{tab:stack}
\small
\begin{tabularx}{\textwidth}{@{}l l X@{}}
\toprule
\textbf{Инструмент} & \textbf{Роль} & \textbf{Обоснование выбора} \\
\midrule
Polars~\cite{polars-docs} & обработка данных & ленивые вычисления и потоковый движок: агрегации и джойны на 48М строк без загрузки в память \\
PyArrow~\cite{pyarrow-docs} & колоночный I/O & батчевое чтение parquet и инкрементальная запись (\texttt{ParquetWriter}) для расчета эмбеддинг-признака \\
PyTorch~\cite{pytorch-docs} & векторные вычисления & расчет эмбеддинг-профилей пользователей и косинусов на GPU/MPS/CPU единым кодом \\
CatBoost~\cite{catboost-docs} & ранжирующая модель & встроенные listwise-лоссы (YetiRank~\cite{yetirank}), нативная поддержка категориальных признаков и пропусков, обучение на GPU \\
scikit-learn~\cite{sklearn-docs} & верификация & эталонная реализация NDCG для сверки собственных метрик в тестах \\
ClearML~\cite{clearml-docs} & трекинг экспериментов & self-hosted сервер, автоматическая фиксация конфигураций, скаляров и артефактов \\
Docker Compose & инфраструктура & локальное развертывание сервера ClearML (API, веб-интерфейс, файловый сервер, MongoDB, Elasticsearch, Redis) \\
pytest & тестирование & фикстуры с цепочкой этапов пайплайна, синтетические данные \\
\bottomrule
\end{tabularx}
\end{table}

\subsection{Выбор системы трекинга экспериментов}

Помимо ClearML рассматривались альтернативы: MLflow~\cite{mlflow-docs}, Weights~\&~Biases~\cite{wandb-docs}, Neptune.ai и TensorBoard.
Сравнение по значимым для проекта критериям приведено в таблице~\ref{tab:tracking-comparison}.

\begin{table}[H]
\centering
\caption{Сравнение систем трекинга экспериментов}
\label{tab:tracking-comparison}
\small
\begin{tabularx}{\textwidth}{@{}l X X@{}}
\toprule
\textbf{Система} & \textbf{Сильные стороны} & \textbf{Ограничения в контексте проекта} \\
\midrule
ClearML & полностью бесплатный self-hosted сервер; автозахват конфигураций, кода, git-состояния и скаляров; артефакты, датасеты, пайплайны и удаленные агенты в одной платформе & образы сервера только под amd64 (решено эмуляцией) \\
MLflow & простой API, стандарт де-факто, self-hosted & нет автозахвата окружения и конфигураций --- все логируется вручную; для артефактов нужно отдельно поднимать хранилище (S3/MinIO); слабее веб-интерфейс сравнения экспериментов \\
Weights \& Biases & богатая визуализация, удобные отчеты & облачный SaaS: данные экспериментов уходят на внешние серверы, что неприемлемо для рабочих данных; self-hosted версия платная \\
Neptune.ai & легковесный API & аналогично W\&B --- облачный сервис с платными тарифами \\
TensorBoard & прост, встроен в экосистему & только визуализация скаляров: нет реестра экспериментов, конфигураций, артефактов и сравнения запусков \\
\bottomrule
\end{tabularx}
\end{table}

Решающими для выбора ClearML стали три фактора.
Во-первых, возможность развернуть \emph{полнофункциональный} сервер локально и бесплатно: эксперименты и артефакты не покидают рабочую машину.
Во-вторых, минимальная инвазивность SDK: одна строка \texttt{Task.init()} на этап автоматически фиксирует конфигурацию, окружение и метрики --- в MLflow сопоставимый уровень фиксации потребовал бы ручного кода в каждом этапе.
В-третьих, запас на развитие: та же платформа предоставляет оркестрацию пайплайнов (ClearML Pipelines), очереди удаленных агентов для обучения на GPU-сервере и подбор гиперпараметров, что закрывает пункты дальнейшего развития проекта без смены инструмента.

\subsection{Развертывание ClearML}

Сервер ClearML развернут локально через Docker Compose и включает семь сервисов: \texttt{apiserver} (порт 8008), \texttt{webserver} (8080), \texttt{fileserver} (8081), MongoDB, Elasticsearch, Redis и служебный \texttt{async\_delete}.
Поскольку рабочая машина построена на архитектуре arm64, а официальные образы ClearML собраны только под amd64, всем сервисам ClearML в \texttt{docker-compose.yml} явно указана платформа \texttt{linux/amd64} --- Docker прозрачно эмулирует ее.

Python-SDK подключается к серверу через конфигурацию \texttt{\textasciitilde/clearml.conf}, создаваемую командой \texttt{clearml-init} по ключам доступа из веб-интерфейса.
Пайплайн спроектирован отказоустойчиво по отношению к трекингу: если сервер недоступен или трекинг выключен в конфигурации, все этапы выполняются штатно, ограничиваясь консольными логами.

\subsection{Архитектура решения}

Решение организовано как последовательность из пяти этапов, каждый из которых читает результаты предыдущего с диска и создает собственный эксперимент в ClearML (таблица~\ref{tab:pipeline-stages}).
Все настройки вынесены в два YAML-конфига: \texttt{configs/pipeline.yaml} (data-этапы) и \texttt{configs/train.yaml} (модель).

\begin{table}[H]
\centering
\caption{Этапы пайплайна}
\label{tab:pipeline-stages}
\small
\begin{tabularx}{\textwidth}{@{}l l X@{}}
\toprule
\textbf{Этап} & \textbf{Модуль} & \textbf{Результат} \\
\midrule
1. prepare & \texttt{src/prepare\_data.py} & \texttt{manifest.yaml}: файлы по категориям, шарды в порядке недель \\
2. preprocess & \texttt{src/preprocess.py} & очищенные события с \texttt{watch\_ratio} и таргетом \texttt{label} \\
3. features & \texttt{src/features.py} & агрегатные и эмбеддинг-признаки по train-окну \\
4. dataset & \texttt{src/build\_dataset.py} & train/val/test parquet + \texttt{dataset\_meta.yaml} \\
5. train & \texttt{src/train.py} & модель \texttt{.cbm}, оффлайн-метрики, важности признаков \\
\bottomrule
\end{tabularx}
\end{table}

Выполненный прогон пайплайна в веб-интерфейсе ClearML показан на рисунке~\ref{fig:clearml-experiments}: каждый этап зафиксирован как отдельный завершенный эксперимент со своим типом (data processing / training).

\begin{figure}[H]
\centering
\includegraphics[width=\textwidth]{clearml_experiments.png}
\caption{Этапы пайплайна в проекте ClearML (веб-интерфейс, вкладка Tasks)}
\label{fig:clearml-experiments}
\end{figure}
